\section{Conclusions}

\myparagraph{Limitations and Future Work.} 
% 
To leverage in-the-wild human actions, we adopt only wrist translation as the bridging action, which limits transfer on tasks that require fine rotational adjustments (Sec.~\ref{exp:sec:failure}). 
% 
We also find the robot struggles to pick up thin objects after co-training, which we attribute mainly to 1) the observation and embodiment gap and 2) the inevitable noise in human actions.
% 
In future work, with larger-scale and more diverse robot actions, we can further narrow the gap between human and robot actions.


\myparagraph{Conclusions.}
% 
In this work, we investigate the feasibility of learning task-specific bi-manual robot skills from human motion data.
% 
We propose a translation-based bridging action representation that is compatible with both robot and human action data. 
% 
To address the potentially missing action components in different data sources, we further introduce an interleaved action sequence representation.
% 
Experimental results show that the proposed translation-only actions can efficiently transfer human manipulation knowledge to robot skill learning and effectively benefit from large-scale pretraining.